\documentclass[11pt]{article}

\usepackage{amsmath, amssymb}
\usepackage{physics}
\usepackage{slashed}
\usepackage{graphicx}

\begin{document}

\section*{$\mu \to e \gamma$}

\subsection*{General calculation}

We consider the radiative charged-lepton flavor violating decay
\begin{equation}
    \mu(p) \to e(p-q) + \gamma(q).
\end{equation}

The amplitude can be written as
\begin{equation}
    \mathcal{M}(\mu \to e\gamma)
    =
    e\, \epsilon^{*\mu}(q)
    \mel{e}{J_\mu^{\rm em}}{\mu}.
\end{equation}

The most general matrix element compatible with Lorentz covariance is
\begin{equation}
    \mel{e}{J_\mu^{\rm em}}{\mu}
    =
    \bar{u}_e(p-q)
    \left[
        i q^\nu \sigma_{\mu\nu} (A + B\gamma_5)
        + \gamma_\mu (C + D\gamma_5)
        + q_\mu (E + F\gamma_5)
    \right]
    u_\mu(p).
\end{equation}

From the Ward identity of electromagnetism,
\begin{equation}
    \partial^\mu J_\mu^{\rm em}=0,
\end{equation}
we have
\begin{equation}
    q^\mu \mel{e}{J_\mu^{\rm em}}{\mu}=0.
\end{equation}

Contracting the general decomposition with $q^\mu$, one obtains
\begin{align}
    q^\mu J_\mu^{\rm em}
    &=
    i q^\mu q^\nu \sigma_{\mu\nu}(A+B\gamma_5)
    + \slashed{q}(C+D\gamma_5)
    + q^2(E+F\gamma_5).
\end{align}

The first term vanishes because $q^\mu q^\nu$ is symmetric while
$\sigma_{\mu\nu}$ is antisymmetric:
\begin{equation}
    q^\mu q^\nu \sigma_{\mu\nu}=0.
\end{equation}

Since the external photon is on shell,
\begin{equation}
    q^2=0,
\end{equation}
the Ward identity reduces to
\begin{equation}
    \bar{u}_e(p-q)\slashed{q}(C+D\gamma_5)u_\mu(p)=0.
\end{equation}

Using
\begin{equation}
    q = p - (p-q),
\end{equation}
and the Dirac equations
\begin{equation}
    \slashed{p} u_\mu(p)=m_\mu u_\mu(p),
    \qquad
    \bar{u}_e(p-q)\slashed{(p-q)}=m_e \bar{u}_e(p-q),
\end{equation}
we find
\begin{align}
    \bar{u}_e(p-q)\slashed{q}(C+D\gamma_5)u_\mu(p)
    &=
    \bar{u}_e(p-q)
    \left[
        \slashed{p} - \slashed{(p-q)}
    \right]
    (C+D\gamma_5)u_\mu(p)
    \notag \\
    &=
    \bar{u}_e(p-q)
    \left[
        m_\mu(C-D\gamma_5)
        -
        m_e(C+D\gamma_5)
    \right]
    u_\mu(p).
\end{align}

Therefore,
\begin{equation}
    m_e(C+D\gamma_5)-m_\mu(C-D\gamma_5)=0.
\end{equation}

In the limit $m_e \ll m_\mu$, this implies
\begin{equation}
    C=D=0.
\end{equation}

Thus the remaining amplitude is
\begin{equation}
    \mathcal{M}(\mu \to e\gamma)
    =
    e\,\epsilon^{*\mu}(q)\,
    \bar{u}_e(p-q)
    \left[
        i q^\nu \sigma_{\mu\nu}(A+B\gamma_5)
    \right]
    u_\mu(p).
\end{equation}

Using the chiral projectors
\begin{equation}
    P_{R,L} = \frac{1\pm \gamma_5}{2},
\end{equation}
we can rewrite
\begin{equation}
    A+B\gamma_5 = A_R P_R + A_L P_L,
\end{equation}
where
\begin{equation}
    A_R = A+B,
    \qquad
    A_L = A-B.
\end{equation}

Therefore,
\begin{equation}
    \mathcal{M}(\mu \to e\gamma)
    =
    e\,\epsilon^{*\mu}(q)\,
    \bar{u}_e(p-q)
    i\sigma_{\mu\nu}q^\nu
    \left(
        A_R P_R + A_L P_L
    \right)
    u_\mu(p).
\end{equation}

After summing over final spins and photon polarizations, and averaging over the initial muon spin, one obtains
\begin{equation}
    \overline{|\mathcal{M}|^2}
    =
    m_\mu^4 e^2
    \left(
        |A_R|^2 + |A_L|^2
    \right),
\end{equation}
where we have used $m_\mu \gg m_e$.

The decay width is then
\begin{equation}
    \Gamma(\mu \to e\gamma)
    =
    \frac{1}{16\pi m_\mu}
    \overline{|\mathcal{M}|^2}
    =
    \frac{m_\mu^3 e^2}{16\pi}
    \left(
        |A_R|^2 + |A_L|^2
    \right).
\end{equation}

\subsection*{SM calculation}

We now compute the contribution in the Standard Model with massive neutrinos. Working in $R_\xi$ gauge and neglecting Goldstone bosons, the relevant diagram contains a $W$ boson and a neutrino running in the loop.

The one-loop amplitude has the schematic form
\begin{align}
    i\mathcal{M}^{(1)}
    &=
    \int \frac{d^4k}{(2\pi)^4}\,
    \bar{u}_e(p-q)
    \sum_k
    \frac{g}{\sqrt{2}} U^*_{ek}
    \gamma^\alpha
    \frac{1-\gamma_5}{2}
    \frac{i(\slashed{p}-\slashed{k}+m_{\nu_k})}
    {(p-k)^2-m_{\nu_k}^2}
    \frac{g}{\sqrt{2}} U_{\mu k}
    \gamma^\beta
    \frac{1-\gamma_5}{2}
    u_\mu(p)
    \notag \\
    &\quad \times
    \frac{-i}{k^2-m_W^2}
    \left(
        g_{\alpha\rho}
        -
        (1-\xi)\frac{k_\alpha k_\rho}{k^2-\xi m_W^2}
    \right)
    \left[
        -ie\,\Gamma^{\rho\sigma\mu}_{WW\gamma}
    \right]
    \notag \\
    &\quad \times
    \frac{-i}{(k+q)^2-m_W^2}
    \left(
        g_{\sigma\beta}
        -
        (1-\xi)\frac{(k+q)_\sigma (k+q)_\beta}
        {(k+q)^2-\xi m_W^2}
    \right)
    \epsilon^*_\mu(q).
\end{align}

The three-gauge-boson vertex is
\begin{equation}
    \Gamma_{\rho\sigma\mu}^{WW\gamma}
    =
    (-q-k)_\sigma g_{\rho\mu}
    +
    (2k+q)_\mu g_{\rho\sigma}
    -
    (2q+k)_\rho g_{\sigma\mu}.
\end{equation}

The fermionic propagator can be expanded as
\begin{equation}
    \sum_k
    \frac{U_{\mu k}U^*_{ek}}
    {(p+k)^2-m_{\nu_k}^2}
    =
    \sum_k
    \left[
        \frac{U_{\mu k}U^*_{ek}}{(p+k)^2}
        +
        \frac{U_{\mu k}U^*_{ek}m_{\nu_k}^2}{(p+k)^4}
        + \cdots
    \right].
\end{equation}

Using unitarity of the PMNS matrix,
\begin{equation}
    UU^\dagger = \mathbb{1},
\end{equation}
we have
\begin{equation}
    \sum_k U_{\mu k}U^*_{ek} = \delta_{\mu e}=0.
\end{equation}

Therefore the leading term vanishes, and the first surviving contribution is proportional to
\begin{equation}
    \sum_k U_{\mu k}U^*_{ek}m_{\nu_k}^2.
\end{equation}

Using the freedom to subtract a common mass, this can also be written as
\begin{align}
    \sum_k U_{\mu k}U^*_{ek}m_{\nu_k}^2
    &=
    -m_1^2
    \left(
        U_{\mu 2}U^*_{e2}
        +
        U_{\mu 3}U^*_{e3}
    \right)
    +
    U_{\mu 2}U^*_{e2}m_2^2
    +
    U_{\mu 3}U^*_{e3}m_3^2
    \notag \\
    &=
    U_{\mu 2}U^*_{e2}\Delta m_{21}^2
    +
    U_{\mu 3}U^*_{e3}\Delta m_{31}^2.
\end{align}

This means that differences of neutrino masses are responsible for the amount of charged lepton flavor violation generated in the SM. In the right-handed amplitude, however, the suppression is even stronger because of the chirality structure.

The loop calculation gives
\begin{equation}
    A^{(1)}_R
    =
    \frac{g^2 e}{32\pi^2}
    \frac{m_\mu}{m_W^4}
    \sum_{k=2,3}
    U_{\mu k}U^*_{ek}m_{\nu_k}^2
    \left[
        -\frac{1}{3}
        +
        \frac{3}{s-1}
        -
        \frac{1}{(s-1)^2}
        \left(
            \frac{2s\ln s}{s-1}
            -1
        \right)
    \right].
\end{equation}

Here $s$ denotes the gauge parameter associated with the $R_\xi$ gauge calculation.

Taking the limit $s \to \infty$, the loop function tends to
\begin{equation}
    A_R^{(1)} \to A_R.
\end{equation}

Thus the decay rate behaves parametrically as
\begin{equation}
    \Gamma(\mu \to e\gamma)
    \propto
    \left|
        \sum_{k=2,3}
        U_{\mu k}U^*_{ek}
        \frac{\Delta m_{k1}^2}{m_W^2}
    \right|^2
    \sim 10^{-44}.
\end{equation}

\subsection*{EFT calculation}

At the effective-field-theory level, the leading contribution to $\mu \to e\gamma$ comes from dimension-six dipole operators. The relevant terms in the effective Lagrangian are
\begin{equation}
    \mathcal{L}
    \supset
    \frac{C^{e\gamma}_{e\mu}}{\Lambda^2}
    \frac{v}{\sqrt{2}}\,
    \bar{e}\sigma_{\mu\nu}P_R\mu\,F^{\mu\nu}
    +
    \frac{C^{e\gamma}_{\mu e}}{\Lambda^2}
    \frac{v}{\sqrt{2}}\,
    \bar{\mu}\sigma_{\mu\nu}P_R e\,F^{\mu\nu}
    + \text{h.c.}
\end{equation}

Equivalently, after taking the hermitian conjugate of the second term, the amplitude can be written as
\begin{equation}
    i\mathcal{M}
    =
    \bar{u}_e(p-q)
    \frac{v}{\sqrt{2}\Lambda^2}
    \sigma_{\mu\nu}q^\mu
    \left[
        C^{e\gamma}_{e\mu}P_R
        +
        C^{e\gamma *}_{\mu e}P_L
    \right]
    u_\mu(p)\,
    \epsilon^{*\nu}(q).
\end{equation}

Matching this amplitude with the general parametrization,
\begin{equation}
    \mathcal{M}(\mu \to e\gamma)
    =
    e\,\epsilon^{*\mu}
    \bar{u}_e i\sigma_{\mu\nu}q^\nu
    \left(
        A_RP_R+A_LP_L
    \right)
    u_\mu,
\end{equation}
we identify
\begin{equation}
    C^{e\gamma}_{e\mu}
    =
    A_R \frac{\sqrt{2}e\Lambda^2}{v},
    \qquad
    C^{e\gamma *}_{\mu e}
    =
    A_L \frac{\sqrt{2}e\Lambda^2}{v}.
\end{equation}

Therefore,
\begin{equation}
    A_R =
    \frac{v}{\sqrt{2}e\Lambda^2}
    C^{e\gamma}_{e\mu},
    \qquad
    A_L =
    \frac{v}{\sqrt{2}e\Lambda^2}
    C^{e\gamma *}_{\mu e}.
\end{equation}

Substituting into the decay width,
\begin{equation}
    \Gamma(\mu \to e\gamma)
    =
    \frac{m_\mu^3 e^2}{16\pi}
    \left(
        |A_R|^2+|A_L|^2
    \right),
\end{equation}
we obtain
\begin{equation}
    \Gamma(\mu \to e\gamma)
    =
    \frac{m_\mu^3 v^2}{32\pi\Lambda^4}
    \left(
        |C^{e\gamma}_{e\mu}|^2
        +
        |C^{e\gamma}_{\mu e}|^2
    \right).
\end{equation}

% In your handwritten note, the final expression appears with an 8\pi
% in the denominator. This depends on the normalization convention for
% the dipole operator and for F_{\mu\nu}. Check whether a factor of 2
% from F_{\mu\nu} = \partial_\mu A_\nu - \partial_\nu A_\mu was included
% in the Feynman rule.

In this EFT description, the suppression does not come from neutrino masses. Instead, it comes from the scale $\Lambda$ of new physics and from the Wilson coefficients $C^{e\gamma}_{e\mu}$ and $C^{e\gamma}_{\mu e}$, that once are abtrary we set that $C\tilde\ \mathcal(O)(1)$.

\end{document}